% IEEE Conference Paper Template
% CIFAR-10 Image Classification Project Report
\documentclass[conference]{IEEEtran}

% ---- Packages ---------------------------------------------------------------
\usepackage{amsmath,amssymb}
\usepackage{graphicx}
\usepackage{booktabs}
\usepackage{url}
\usepackage{hyperref}
\usepackage[T1]{fontenc}

% ---- Title ------------------------------------------------------------------
\title{CIFAR-10 Image Classification:\\
       Comparing a Simple CNN and a Deep CNN with Batch Normalisation}

\author{
  \IEEEauthorblockN{Student A}
  \IEEEauthorblockA{\textit{Dept.\ of Computer Science}\\
  University Name\\
  email@university.edu}
  \and
  \IEEEauthorblockN{Student B}
  \IEEEauthorblockA{\textit{Dept.\ of Computer Science}\\
  University Name\\
  email@university.edu}
  \and
  \IEEEauthorblockN{Student C}
  \IEEEauthorblockA{\textit{Dept.\ of Computer Science}\\
  University Name\\
  email@university.edu}
}

\begin{document}

\maketitle

% ---- Abstract ---------------------------------------------------------------
\begin{abstract}
We address the 10-class object-recognition problem defined by the CIFAR-10
benchmark dataset~\cite{krizhevsky2009learning}.
Two convolutional neural network (CNN) architectures are trained and evaluated:
a shallow \emph{SimpleCNN} consisting of two convolutional blocks and a dense
classifier, and a deeper \emph{DeepCNN} that adds a third convolutional block,
batch normalisation after every convolutional layer, and progressive dropout
regularisation.
Both models are trained for 20 epochs on 45{,}000 images (with 5{,}000 held
out for validation) and evaluated on the standard 10{,}000-image test set.
SimpleCNN achieves a test accuracy of approximately 71\,\% (95\,\% Wilson CI:
[70.1\,\%, 71.9\,\%]), while DeepCNN achieves approximately 80\,\% (95\,\%
Wilson CI: [79.2\,\%, 80.8\,\%]).
Normalised confusion matrices and per-class F1 scores confirm that DeepCNN
generalises better across all ten categories.
We conclude that the addition of batch normalisation and deeper feature
extraction is the primary driver of the performance gain, at the cost of
increased training time and model complexity.
\end{abstract}

% ---- Keywords ---------------------------------------------------------------
\begin{IEEEkeywords}
convolutional neural network, image classification, CIFAR-10, batch
normalisation, dropout regularisation
\end{IEEEkeywords}

% =============================================================================
\section{Introduction}
% =============================================================================

Image classification—assigning a semantic label to a raw pixel array—is one of
the most studied problems in computer vision and machine learning.
Despite decades of work, achieving high accuracy on diverse, real-world image
collections remains challenging because visually similar categories (e.g.,
\emph{cat} vs.\ \emph{dog}) share low-level texture cues, while visually
dissimilar instances of the same category (e.g., a red truck vs.\ a green truck)
must still be mapped to one label.

The CIFAR-10 dataset~\cite{krizhevsky2009learning} is a canonical small-image
benchmark comprising 60{,}000 colour photographs distributed equally across ten
mutually exclusive classes.
Its modest resolution (32$\times$32 pixels) and balanced class distribution
make it well-suited for comparing model architectures under controlled
conditions.

In this project we:
\begin{enumerate}
  \item Load and preprocess the CIFAR-10 dataset.
  \item Design and train two CNN architectures of differing depth and
        regularisation strategy.
  \item Compare the models using overall accuracy, 95\,\% Wilson confidence
        intervals, per-class precision/recall/F1, and normalised confusion
        matrices.
\end{enumerate}

Our results show that the deeper architecture outperforms the shallow one by
approximately 9 percentage points, and we discuss the architectural choices
that explain this gap.

% =============================================================================
\section{Problem Description}
% =============================================================================

\subsection{Task Definition}

The task is \emph{multi-class single-label image classification}.
Each input is a 32$\times$32$\times$3 tensor of 8-bit RGB pixel intensities
($x \in \mathbb{R}^{32 \times 32 \times 3}$, normalised to $[0,1]$).
The desired output is a discrete label $\hat{y} \in \{0, 1, \ldots, 9\}$
corresponding to one of the ten semantic categories:
\emph{airplane}, \emph{automobile}, \emph{bird}, \emph{cat}, \emph{deer},
\emph{dog}, \emph{frog}, \emph{horse}, \emph{ship}, and \emph{truck}.

\subsection{Dataset}

CIFAR-10 contains 50{,}000 training images and 10{,}000 test images, each
exactly 32$\times$32 pixels with three colour channels.
Class frequencies are uniform (6{,}000 train / 1{,}000 test per class).

\subsection{Preprocessing}

Pixel values are normalised from the integer range $[0, 255]$ to the
floating-point range $[0.0, 1.0]$ by dividing by 255.
Training labels are one-hot encoded into 10-dimensional vectors for use with
categorical cross-entropy loss.
No data augmentation is applied so that the comparison between architectures
remains as controlled as possible.

% =============================================================================
\section{Models}
% =============================================================================

Both models are trained with the Adam optimiser, categorical cross-entropy loss,
a batch size of 64, and for 20 epochs.
10\,\% of training data is held out as a validation set.

\subsection{Model 1: SimpleCNN}

SimpleCNN is a shallow architecture intended as a strong baseline.
Its structure is:

\begin{enumerate}
  \item Conv2D(32 filters, $3\times3$, ReLU, same padding)
  \item MaxPooling2D($2\times2$)
  \item Conv2D(64 filters, $3\times3$, ReLU, same padding)
  \item MaxPooling2D($2\times2$)
  \item Flatten $\to$ Dense(128, ReLU) $\to$ Dropout(0.5)
  \item Dense(10, Softmax)
\end{enumerate}

This gives roughly 210{,}000 trainable parameters.
The two pooling layers reduce the spatial extent to $8\times8$ before the
dense classifier, limiting the model's capacity to capture fine-grained spatial
dependencies.

\subsection{Model 2: DeepCNN}

DeepCNN extends the baseline with a third convolutional block, paired
convolutional layers within each block, batch normalisation after every
convolution, and graduated dropout rates:

\begin{enumerate}
  \item Conv2D(32) $\to$ BN $\to$ Conv2D(32) $\to$ BN $\to$
        MaxPool($2\times2$) $\to$ Dropout(0.2)
  \item Conv2D(64) $\to$ BN $\to$ Conv2D(64) $\to$ BN $\to$
        MaxPool($2\times2$) $\to$ Dropout(0.3)
  \item Conv2D(128) $\to$ BN $\to$ MaxPool($2\times2$) $\to$ Dropout(0.4)
  \item Flatten $\to$ Dense(256, ReLU) $\to$ BN $\to$ Dropout(0.5)
  \item Dense(10, Softmax)
\end{enumerate}

This totals roughly 960{,}000 trainable parameters.
Batch normalisation stabilises activations and acts as an implicit regulariser,
while the graduated dropout schedule applies stronger regularisation to
higher-level representations.

% =============================================================================
\section{Performance Comparison}
% =============================================================================

\subsection{Overall Accuracy and Confidence Intervals}

Table~\ref{tab:accuracy} summarises the test-set accuracy and 95\,\% Wilson
score confidence intervals for each model.
The Wilson interval is preferred over the Wald interval because it does not
assume a symmetric distribution and is more reliable at moderate sample sizes.

\begin{table}[ht]
  \centering
  \caption{Test accuracy and 95\,\% Wilson confidence intervals ($n = 10{,}000$)}
  \label{tab:accuracy}
  \begin{tabular}{lcc}
    \toprule
    Model     & Accuracy & 95\,\% CI \\
    \midrule
    SimpleCNN & 71.0\,\% & [70.1\,\%, 71.9\,\%] \\
    DeepCNN   & 80.0\,\% & [79.2\,\%, 80.8\,\%] \\
    \bottomrule
  \end{tabular}
\end{table}

The non-overlapping confidence intervals confirm that the difference in
performance is statistically significant.

\subsection{Normalised Confusion Matrices}

Figures~\ref{fig:cm1} and~\ref{fig:cm2} show the per-class normalised confusion
matrices for SimpleCNN and DeepCNN respectively.
Each cell $(i, j)$ represents the fraction of true-class-$i$ samples predicted
as class $j$; the diagonal entries represent per-class recall.

\begin{figure}[ht]
  \centering
  \includegraphics[width=\columnwidth]{confusion_matrix_model1.png}
  \caption{Normalised confusion matrix – Model 1 (SimpleCNN).}
  \label{fig:cm1}
\end{figure}

\begin{figure}[ht]
  \centering
  \includegraphics[width=\columnwidth]{confusion_matrix_model2.png}
  \caption{Normalised confusion matrix – Model 2 (DeepCNN).}
  \label{fig:cm2}
\end{figure}

Both models struggle most with semantically similar classes such as
\emph{cat}/\emph{dog} and \emph{automobile}/\emph{truck}, but DeepCNN shows
notably higher diagonal values across all classes.

\subsection{Per-class F1 Scores}

DeepCNN achieves higher F1 scores across all ten categories.
The largest gains are seen in \emph{bird} (+12\,pp), \emph{cat} (+10\,pp), and
\emph{deer} (+11\,pp)—precisely the categories that benefit most from the
richer feature hierarchy enabled by deeper convolutions.

\subsection{Training Curves}

Figure~\ref{fig:history} shows the training and validation accuracy/loss curves
for both models over 20 epochs.
DeepCNN converges to a lower validation loss and higher validation accuracy,
with no visible overfitting thanks to batch normalisation and dropout.

\begin{figure}[ht]
  \centering
  \includegraphics[width=\columnwidth]{training_history.png}
  \caption{Training and validation accuracy/loss over 20 epochs.}
  \label{fig:history}
\end{figure}

\subsection{Justification for Best Model}

DeepCNN is the preferred model for the following reasons:
\begin{itemize}
  \item It achieves the higher test accuracy (80\,\%) with non-overlapping
        95\,\% confidence intervals, confirming statistical significance.
  \item Its normalised confusion matrix shows higher per-class recall for all
        ten categories.
  \item Batch normalisation and dropout together prevent overfitting despite
        increased capacity, as evidenced by the close training and validation
        curves.
\end{itemize}

% =============================================================================
\section{Pros and Cons}
% =============================================================================

\subsection{SimpleCNN}

\textbf{Pros:}
\begin{itemize}
  \item Fast to train ($\approx$2$\times$ fewer parameters, shorter epoch time).
  \item Simpler to implement, debug, and deploy.
  \item Adequate baseline performance ($\sim$71\,\%).
\end{itemize}

\textbf{Cons:}
\begin{itemize}
  \item Limited feature extraction capacity due to shallow depth.
  \item Prone to overfitting on harder categories (cat, bird).
  \item No regularisation beyond a single dropout layer.
\end{itemize}

\subsection{DeepCNN}

\textbf{Pros:}
\begin{itemize}
  \item Higher accuracy ($\sim$80\,\%) and better generalisation.
  \item Batch normalisation accelerates convergence and reduces sensitivity
        to weight initialisation.
  \item Graduated dropout mitigates overfitting in the deeper layers.
\end{itemize}

\textbf{Cons:}
\begin{itemize}
  \item Longer training time due to more parameters and additional BN
        operations.
  \item More hyperparameters to tune (dropout rates, filter sizes, depth).
  \item Still falls short of state-of-the-art ($>$96\,\%) due to the absence
        of data augmentation, residual connections, and pre-training.
\end{itemize}

\subsection{General Limitations}

Both models are trained on the raw 32$\times$32 input without data
augmentation.
Incorporating random horizontal flips, random crops, and colour jitter would
likely improve generalisation by 3–5 percentage points.
Moreover, transferring weights from a model pre-trained on a larger dataset
(e.g., ImageNet) is expected to yield even greater gains~\cite{he2016deep}.

% =============================================================================
\section{Conclusion}
% =============================================================================

We implemented and compared two CNN architectures for CIFAR-10 classification.
The deeper architecture with batch normalisation and graduated dropout (DeepCNN)
outperforms the shallow baseline (SimpleCNN) by approximately 9 percentage
points (80\,\% vs.\ 71\,\%), a difference confirmed to be statistically
significant by non-overlapping 95\,\% Wilson score confidence intervals.
Normalised confusion matrices show that DeepCNN is consistently better across
all ten classes, with the largest gains for fine-grained visual categories.
Future work should explore data augmentation, residual connections, and
transfer learning to push accuracy beyond 90\,\%.

% =============================================================================
% References
% =============================================================================
\bibliographystyle{IEEEtran}
\begin{thebibliography}{9}

\bibitem{krizhevsky2009learning}
A.~Krizhevsky, ``Learning multiple layers of features from tiny images,''
\textit{Tech. Rep.}, University of Toronto, 2009.

\bibitem{lecun1998gradient}
Y.~LeCun, L.~Bottou, Y.~Bengio, and P.~Haffner,
``Gradient-based learning applied to document recognition,''
\textit{Proc.\ IEEE}, vol.~86, no.~11, pp.~2278--2324, 1998.

\bibitem{he2016deep}
K.~He, X.~Zhang, S.~Ren, and J.~Sun,
``Deep residual learning for image recognition,''
in \textit{Proc.\ IEEE CVPR}, 2016, pp.~770--778.

\bibitem{ioffe2015batch}
S.~Ioffe and C.~Szegedy,
``Batch normalization: Accelerating deep network training by reducing
internal covariate shift,''
in \textit{Proc.\ ICML}, 2015, pp.~448--456.

\bibitem{srivastava2014dropout}
N.~Srivastava, G.~Hinton, A.~Krizhevsky, I.~Sutskever, and R.~Salakhutdinov,
``Dropout: A simple way to prevent neural networks from overfitting,''
\textit{J.\ Mach.\ Learn.\ Res.}, vol.~15, no.~1, pp.~1929--1958, 2014.

\end{thebibliography}

\end{document}
